\section{Introduction}
Cold-spray particle impact involves large plastic deformation, severe strain-rate amplification near the contact interface, and rapid conversion of plastic work into heat \cite{JohnsonCook1983,TaylorQuinney1934}. In this regime, a constitutive law must reproduce at least four coupled effects: strain hardening, low-to-moderate strain-rate strengthening, thermal softening, and an additional high-rate resistance once the plastic strain rate exceeds a transition scale. A conventional Johnson--Cook (JC) law provides a practical starting point for the first three of these effects, but its standard logarithmic rate term is generally too weak to reproduce the sharp rise of flow stress observed in the extreme-rate regime.

Recent impact and indentation studies reinforce this point. High-velocity microparticle impact experiments and simulations show that rebound, flattening, jetting, and adhesion do not evolve smoothly under a single weak rate-sensitive law; instead, they exhibit a marked transition as the impact severity increases \cite{Tiamiyu2021,Zhang2025Peridynamics}. In addition, indentation-based ultra-high-strain-rate reconstructions indicate that representative flow strength rises strongly in the range from about $10^5$ to $10^6$ s$^{-1}$ and above, which is not well captured by the standard JC logarithmic rate term alone \cite{Hassani2019,Chen2026Indentation}. These observations provide the practical constitutive motivation for introducing a separate high-rate strengthening branch.

The motivation for the present work is therefore twofold. First, a low-rate JC branch is retained because it remains a compact and interpretable engineering description of rate- and temperature-dependent plastic flow. Second, a separate high-rate branch is introduced to represent the additional drag-dominant strengthening associated with fast dislocation motion in a thermally activated crystal lattice. This additional branch is not intended as a full dislocation-dynamics model. Instead, it is a physics-informed continuum closure designed for impact simulations in which a tractable constitutive form is needed, but where the transition to drag-controlled behavior must still be tied to physically interpretable scales such as the Burgers vector, the shear modulus, and a characteristic transverse sound speed \cite{Blaschke2019PhononWind,Orowan1940}.

A central design choice in this work is that the effective transition strain rate is treated as a material-dependent constitutive threshold, not as a direct function of the externally prescribed impact velocity. The impact velocity determines the actually attained strain-rate field in the particle and substrate, whereas the constitutive transition rate determines when the drag-dominant strengthening mechanism begins to contribute. This distinction proved important during model development and is preserved explicitly in the final formulation.

The paper is organized as follows. Section~2 introduces the phonon-drag-motivated constitutive closure and its coupling to the JC baseline and the Taylor--Quinney heat-generation term. Section~3 summarizes the numerical implementation used in the peridynamic impact solver. The emphasis throughout is on preserving a transparent link between the continuum constitutive form and the material scales used to interpret the transition to the high-rate regime.
